% ShivyCX rpython cross-runtime benchmark report.
%
% This is the static shell. The data-driven body (tables + figures) lives in
% report_body.tex, which benchmarks/plot_rpython.py regenerates from the latest
% results/rpython_results.json. Build it with `make benchmarks-report`, which
% runs the harness, renders the figures, and invokes pdflatex in the figure
% output directory (so report_body.tex and the *.pdf figures sit alongside).
\documentclass[11pt]{article}
\usepackage[margin=1in]{geometry}
\usepackage{graphicx}
\usepackage{booktabs}
\usepackage{array}
\usepackage{xcolor}
\usepackage{hyperref}
\hypersetup{colorlinks=true, linkcolor=blue, urlcolor=blue}

\title{ShivyCX benchmarks: minipy and the cross-runtime suite\\
\large A \texttt{py2c}-compiled Python interpreter, plus
CPython vs.\ PyPy3 vs.\ py2c$+$gcc vs.\ the self-hosted compiler}
\date{\today}

\begin{document}
\maketitle

\input{minipy_body.tex}

\section{What is measured}

Each benchmark in \texttt{benchmarks/rpython/} is a single \emph{pure-Python}
program. The same source is executed by four backends, which must all agree on
the program's exit code (a differential-correctness check):

\begin{itemize}
  \item \textbf{CPython} --- the program run directly under \texttt{python3}.
  \item \textbf{PyPy3} --- the same program under the PyPy3 JIT (omitted if PyPy3
        is not installed).
  \item \textbf{py2c$+$gcc} --- the program transpiled to C by
        \texttt{tools/py2c.py} and compiled with \texttt{gcc -O2}.
  \item \textbf{self-hosted} --- the \emph{same} transpiled C compiled by
        \texttt{shivyc\_native}, the ShivyCX compiler built from its own
        sources. This is the headline column: ShivyCX compiling the C that its
        own \texttt{.py}$\rightarrow$C translator produced.
\end{itemize}

For every backend three quantities are recorded:

\begin{itemize}
  \item \textbf{Runtime} --- best of five wall-clock measurements of the program
        proper.
  \item \textbf{Peak memory} --- the exact maximum resident set size of the
        process, read from the child's \texttt{rusage} (\texttt{ru\_maxrss}) by
        a $\sim$1.5\,MB C probe, so the harness's own footprint does not inflate
        it.
  \item \textbf{Compile time} --- for the two compiled backends, the wall-clock
        time the C compiler (gcc or \texttt{shivyc\_native}) spent on the py2c
        output. The interpreters have no compile step.
\end{itemize}

ShivyCX emits unoptimised, \texttt{-O0}-class code, so the self-hosted runtime is
expected to trail \texttt{gcc -O2}; the interesting questions are how it compares
to the \emph{interpreters} (it beats them on compute-bound, allocation-light
programs and trails them where boxed-list element access dominates) and how its
compile speed and the resulting binary's memory footprint compare to gcc.

\input{report_body.tex}

\section{Reproducing}

\begin{verbatim}
make benchmarks-report      # harness -> figures -> this PDF
\end{verbatim}

\noindent
The harness builds \texttt{shivyc\_native} once (or reuses one supplied via
\texttt{\$SHIVYC\_NATIVE}); results land in
\texttt{benchmarks/results/rpython\_results.json} and the figures in
\texttt{\$BENCH\_PLOT\_DIR} (default \texttt{/tmp/shivyc\_benchmarks}).

\end{document}
